\begin{textsample}{POS Dim 1 – vem\_ed\_34 – Score 88.00 – vem\_ed\_34/t185.txt}  \label{ex:f1_pos_006}
Artigo de Opinião Cultivando PCIs duradouros Os Projetos Colaborativos Internacionais (PCIs) \textbf{são} atividades de \textbf{extensão} \textbf{institucional} das Fatecs, \textbf{que} oferecem oportunidades de Internacionalização em Casa para os estudantes, desenvolvendo \textbf{competências} \textbf{linguísticas}, \textbf{socioemocionais}, digitais e interculturais, imprescindíveis para a atuação no \textbf{mercado} de trabalho.
A equipe de Apoio à Internacionalização ao Ensino Superior (AIES) da Divisão de Extensão e Pesquisa no Ensino Superior (Depes) da Coordenadoria Geral de Ensino Superior de Graduação (CGESG) do CPS se reuniu no \textbf{início} de abril para uma \textbf{reflexão} coletiva a \textbf{partir} de duas \textbf{questões} norteadoras: 1.
Quais os \textbf{fatores} de sucesso dos PCIs? 2.
Como planejar, executar e avaliar um PCI para \textbf{que} se \textbf{torne} uma parceria estável e duradoura?
De \textbf{acordo} com as discussões da equipe, um dos principais \textbf{fatores} de sucesso \textbf{é} a sintonia entre \textbf{todos} os participantes.
Sintonia vai além da afinidade entre docentes: envolve alinhamento de \textbf{expectativas}, objetivos de aprendizagem convergentes, \textbf{compreensão} \textbf{mútua} dos \textbf{contextos} \textbf{institucionais} e disposição para o \textbf{diálogo} intercultural.
Quando professores, estudantes e equipes de \textbf{coordenação} compartilham propósitos claros, o projeto tende a fluir de \textbf{forma} mais orgânica.
Outro \textbf{elemento} central \textbf{é} a autonomia, tanto docente \textbf{quanto} discente.
PCIs eficazes \textbf{são} aqueles \textbf{que} \textbf{reconhecem} os professores como autores do projeto e os estudantes como sujeitos \textbf{ativos} do \textbf{processo} de aprendizagem. \textbf{Cada} participante \textbf{é} corresponsável pelo sucesso coletivo do PCI.
Esse \textbf{ponto} se \textbf{conecta} com outro aspecto \textbf{essencial}, a colaboração efetiva.
Não funciona com a fórmula “\textbf{cada} grupo faz sua \textbf{parte}, depois apresenta ao outro e \textbf{compara} os resultados no final do projeto”.
Os grupos de estudantes \textbf{são} mistos (Fatec + instituição internacional) e as atividades devem \textbf{ser} \textbf{desenvolvidas} em conjunto.
Os PCIs, dentro da abordagem COIL (Collaborative Online International Learning), pressupõem \textbf{interação} \textbf{significativa} entre estudantes de diferentes países, mediada por \textbf{tarefas} \textbf{colaborativas}, negociação de sentidos e \textbf{construção} conjunta de \textbf{conhecimento}. \textbf{É} nesse percurso – e não apenas no \textbf{produto} – \textbf{que} se desenvolvem \textbf{competências} acadêmicas, \textbf{profissionais} e \textbf{socioemocionais}.
Somam-se a esses \textbf{fatores} a gestão do tempo e o planejamento realista. \textbf{Diferenças} de calendário acadêmico, fusos \textbf{horários} e \textbf{rotinas} \textbf{institucionais}, incluindo o perfil dos fatecanos (trabalhadores \textbf{que} estudam), \textbf{tornam} o tempo um dos principais desafios dos PCIs.
Pesquisas de percepção realizadas com discentes participantes de PCIs em 2021 e 2022 e \textbf{publicadas} em edições \textbf{anteriores} de VEm (\textbf{números} 8, 12 e 14) mostram \textbf{que} “gerenciar o tempo” está entre os principais desafios encontrados.
PCIs \textbf{bem-sucedidos} respeitam essas limitações, \textbf{estabelecem} \textbf{cronogramas} \textbf{possíveis} e preservam margens de adaptação ao \textbf{longo} do caminho.
Antes de cultivar, \textbf{é} preciso semear O planejamento de um PCI \textbf{começa} muito antes do contato entre estudantes.
Ele se \textbf{inicia} no chamado matchmaking: a identificação de parceiros cujas \textbf{expectativas}, áreas de \textbf{conhecimento} e objetivos de aprendizagem dialoguem entre si.
Projetos podem \textbf{ser} intra, inter ou transdisciplinares.
O \textbf{que} \textbf{importa} \textbf{é} ter clareza sobre o \textbf{que} se espera \textbf{aprender}, ensinar e \textbf{construir} colaborativamente.
A localização de parceiros, o design e acompanhamento dos PCIs ocorrem com a mentoria da equipe da AIES / Depes / CGESG.
Um bom design de projeto funciona como uma âncora em mares instáveis.
Embora a imprevisibilidade \textbf{seja} inerente aos PCIs, especialmente em \textbf{contextos} interculturais, \textbf{referências} claras ajudam docentes e estudantes a se orientarem ao \textbf{longo} do percurso.
Planejar, nesse sentido, não \textbf{é} engessar, mas \textbf{criar} condições para \textbf{que} o projeto se sustente mesmo diante de ajustes necessários.
Traçar um mapa para transitar no território sem se perder, apesar das \textbf{mudanças} de \textbf{rota}.
Na execução, a \textbf{flexibilidade} assume \textbf{papel} central.
Os PCIs / COIL exigem \textbf{capacidade} de renegociação constante, escuta \textbf{ativa} e abertura para o erro como \textbf{parte} do \textbf{processo} formativo.
A experiência mostra \textbf{que} resultados \textbf{significativos} não estão necessariamente associados a \textbf{produtos} / entregáveis complexos, mas a percursos coerentes, bem mediados e \textbf{relevantes} para os estudantes.
A avaliação, por sua vez, precisa \textbf{ser} compreendida em múltiplas camadas.
Avalia-se o \textbf{desenvolvimento} do estudante, individual e coletivamente; o próprio projeto; e a parceria \textbf{estabelecida}.
Instrumentos avaliativos \textbf{compartilhados}, \textbf{reflexões} conjuntas entre docentes parceiros e a \textbf{troca} intercultural sobre critérios e percepções ampliam a qualidade do \textbf{processo} e \textbf{fortalecem} os vínculos \textbf{institucionais}.
Para cultivar PCIs longevos, \textbf{é} preciso \textbf{contar} com suporte \textbf{institucional} e reciprocidade entre as \textbf{coordenações} (Centro Paula Souza e instituições internacionais).
Parcerias duradouras dependem não apenas do entusiasmo individual de professores, mas de vínculos \textbf{institucionais} capazes de garantir continuidade, memória do \textbf{processo} e \textbf{apoio} diante dos desafios inevitáveis.
Além disso, a prática da autocrítica da equipe se mostra \textbf{essencial}.
Ouvir parceiros externos, \textbf{analisar} dados das Pesquisas de Percepção realizadas com alunos e professores e revisitar experiências \textbf{anteriores} \textbf{permite} \textbf{identificar} \textbf{pontos} fortes, fragilidades e caminhos de aprimoramento \textbf{contínuo}.
Parcerias duradouras se constroem justamente nessa disposição para pesquisar o \textbf{processo} e \textbf{aprender} com a própria prática.
Para colher, \textbf{é} preciso podar e adubar À medida \textbf{que} os PCIs se institucionalizam, cresce o risco de \textbf{que} \textbf{processos} burocráticos se sobreponham à dimensão criativa e cultural dos projetos.
Manter o equilíbrio entre organização e abertura \textbf{é} um desafio permanente.
Priorizar qualidade e \textbf{impacto} formativo, em vez de apenas \textbf{números} e \textbf{produtos}, \textbf{é} fundamental para preservar a essência \textbf{colaborativa} desses projetos.
A experiência \textbf{acumulada} pela equipe de \textbf{coordenação} dos PCIs aponta \textbf{que} o sucesso dessas iniciativas não está em fórmulas prontas, mas na combinação sensível entre planejamento cuidadoso e \textbf{flexibilidade}, colaboração genuína e \textbf{reflexão} \textbf{contínua}.
Para colher resultados duradouros, \textbf{é} preciso \textbf{saber} o tempo de podar e a hora de adubar. \textbf{É} necessário abraçar a complexidade do \textbf{encontro} intercultural, \textbf{reconhecer} o valor do percurso e compreender \textbf{que} parcerias sólidas se cultivam ao \textbf{longo} do tempo.
Mais do \textbf{que} projetos, trata-se de \textbf{construir} relações – e é nelas \textbf{que} reside a \textbf{verdadeira} potência das \textbf{ações} de Internacionalização em Casa proporcionadas pelos PCIs.

% matched lemmas: acordo, acumular, analisar, anterior, apoio, aprender, ativo, ação, bem-sucedido, cada, capacidade, colaborativo, começar, comparar, compartilhar, competência, compreensão, conectar, conhecimento, construir, construção, contar, contexto, contínuo, coordenação, criar, cronograma, desenvolvimento, desenvolvir, diferença, diálogo, elemento, encontro, essencial, estabelecer, expectativa, extensão, fator, flexibilidade, forma, fortalecer, horário, identificar, impacto, importar, iniciar, institucional, interação, início, linguístico, longo, mercado, mudança, mútuo, número, papel, parte, partir, permitir, ponto, possível, processo, produto, profissional, publicar, quanto, que, questão, reconhecer, referência, reflexão, relevante, rota, rotina, saber, ser, significativo, socioemocional, tarefa, todo, tornar, troca, verdadeiro
\end{textsample}
